% Notes on queen endings. See README.md for the format.

% Diagram 138 of the 1500-1599 book
@ 8/6q1/8/5QP1/2k5/5K2/8/8 b - - 0 73 | Qc3+?
73...\K b3! puts the king near the corner diagonally opposite the queening square, where it is
safest from counter-checks (\xref{Section}{queen-and-pawn}).\index{Queen endings!Drawing zone}

% Diagram 139 of the 1500-1599 book, left out
@ 8/8/8/K7/3Q4/2P2qk1/8/8 w - - 2 68 | skip

% Diagram 140 of the 1500-1599 book
@ 8/5k2/8/4Q3/6K1/6P1/7q/8 b - - 29 84 | Kg6?
Black must realize that the pawn ending is lost for him, so he must avoid any chance of a queen
exchange. The text move allows 85.\Q g5+ \K f7 86.\Q h5+, forcing the queens off.\index{Transposing into a pawn ending}

@ 8/8/6k1/4Q3/6K1/6P1/7q/8 w - - 30 85 | Qe6+?
White misses his chance: 85.\Q g5+! \K f7 86.\Q h5+ wins.\index{Transposing into a pawn ending}

% Diagram 141 of the 1500-1599 book
@ K7/P7/1Q6/8/q4k2/8/8/8 b - - 0 75 | Qe4+?
75...\K g3! The black king heads for the corner diagonally opposite the queening square, here h1,
where it is safest from counter-checks.\index{Queen endings!Drawing zone} Checks do not help:
75...\Q e8+? 76.\Q b8+ answers check with check.\index{Queen endings!Counter-check} Please pay
attention: a perpetual check is not a threat here.

@ K7/P7/1Q6/8/4qk2/8/8/8 w - - 1 76 | Qb7?
76.\K b8! was the try that wins.
